\section{Positive score recovery from constrained samples}
\label{cube:transfer-section}

Detection in Section~\ref{cube:section} supplies finite linear models,
but their coefficients can have either sign. The next proposition
returns a positive score against a fixed patch to the parent box.
The comparison weights used in its proof have pre bounds; only the
final conversion of such a weight into a patch may charge later
scales and the heights of bases of the constraining spaces.

\begin{proposition}[Transfer of a fixed scalar score]
\label{cube:conditioned-transfer}
Assume the uniform constrained-sampler inputs of
Proposition~\ref{cube:comparison}.  Fix $\epsilon>0$ and prescribe
pre bounds for the input complexities and for $\rho^{-1},\tau^{-1}$,
where $\rho,\tau>0$.  Include these data in the prescribed finite
collection of sampler requests and choose the later parameters in
the order of Section~\ref{cube:sampler-proof}.
Fix a tuple of centers and sample $\psi:\mathcal T\to U$ by the
normalized law~\eqref{sampling:tilted-law}.  Suppose
$f:U\to[0,1]$, $\lambda\in[0,1]$, and a fixed positive patch formula
$T_*(u,\beta)$ of the prescribed complexity with $d_*$ additional
slots satisfy
\[
 \Pr_\psi\left\{
 \begin{array}{l}
 \text{there exists an allowed slice }B\subseteq\mathcal T
 \text{ such that}\\[2pt]
 \displaystyle
 \E_{t\in B}(f(\psi(t))-\lambda)
       T_*(\psi(t),\beta(\psi(t)))\ge\rho
 \end{array}\right\}\ge\tau.
\]
The inner average is uniform on $B$, which may be chosen separately
for each path; the probability uses the fixed-center law without
conditioning on this event.  Then there is an ambient slice and a
positive patch with at most
\[
                         d_*+\sum_h\dim W_h
\]
slots having positive score against $f-(1-\epsilon)\lambda$.
The output complexity and the logarithm of the reciprocal score may
additionally use the allowed late parameters and the heights of bases
of the $W_h$: both have polynomial bounds in the resulting augmented
logarithmic budget, as does the cost of the output slice.
\end{proposition}

We first define a nonnegative comparison weight for a fixed path.
Modular rank will give its uniform cap and its approximation by
site twists, independently of the number of prime powers in the
comparison law. These bounds let us forecast the finite model
complexity before fixing the ensemble. We then compare the model
atoms and recover the determining slots. The final subsection
checks that this argument and the preceding cube comparison use
one ensemble with one order of parameter choices.

\subsection{A positive comparison weight for a fixed path}

Fix a path and one allowed slice. Write $t=(t_\parallel,t_I)$,
where $t_I$ comprises the inactive blocks and $t_\parallel$
comprises all kernel, continuous, and active blocks. Every
parallel side has length $L$ before taking the slice. At a true
site $u$, use the chart coordinates $y_W(u),z(u),w(u)$ of
Section~\ref{sampling:constrained-section}, and put
\[
 b(u)=(u,w(u),z_{\rm act}(u)).
\]
On the fixed path these are integer polynomials
$b(t)=(\psi(t),w(t),z_{\rm act}(t))$. The spatial coordinates
have degree tag $1$ and the coordinates belonging to layer $h$
have degree tag $h$. Let $m_{\rm out}$ be their total number.

A slice may have different strides $D_i$ in its parallel coordinates.
Put $D=\operatorname{lcm}_iD_i$, so that
$\log D\le\sum_i\log D_i$. Partition each parallel progression
into its nonempty residue classes modulo $D$, retaining their
original probability weights. The inactive coordinates are not
changed. Since $D$ is pre-bounded and the parallel sides are
chosen later, these pieces remain permitted pre-cost slices.
We include these refinements in the prescribed slice family.
There are at most $D^{\dim t_\parallel}$ pieces, a pre-bounded
number. A productive original slice has a refined piece with at
least the same score. Choosing that piece separately on each
productive path loses no productive probability. Thus, for score
transfer, we may work on one refined slice, on which all parallel
congruences have modulus $D$.

There are three parts to the comparison law. The inactive
parameters retain their actual integer values. Parallel
parameters have separate real and residue copies, sampled
independently conditional on that same inactive vector.

For the real part, the independent blocks give a continuous
density with pre cap and Lipschitz bound. For each lift axis,
apply the continuous product-block estimate with the principal
coefficients fixed in their prescribed normalized intervals
away from zero. For the spatial output use at least $n+2$
independent nonsingular kernel blocks. If $A$ is one normalized
kernel matrix, its inverse bound ensures that some coordinate
of $A^{\mathsf T}\xi$ has magnitude at least a pre-large
multiple of $|\xi|$. Integration in that parameter gives decay
$B/(1+|\xi|)$. The independent blocks multiply these bounds,
giving integrable Fourier decay with a first moment in $n$
dimensions. The kernel and lift blocks use disjoint parameters.
Fourier inversion therefore gives a continuous joint density;
we use this version at every output point. The same estimates
hold after the allowed coordinatewise slice restrictions.

Write $G$ for this density of
\[
             (u/H,y_W,z_{\rm act}/K)
\]
obtained from real normalized parallel parameters on the slice.
Use the principal lift polynomials, including their actual
constant shifts, and the kernel contribution to spatial
displacement. These outputs depend only on $t_\parallel$.
The small nonprincipal lift terms and nonkernel spatial
displacement are omitted in this real density. Their smallness
will be used only when comparing smooth test values.

For the residue part, choose a finite set of distinct primes
and one positive power of each, and let $N$ be their product.
Given the actual $t_I$,
sample $t_\parallel\bmod N$ uniformly subject to the slice
congruences, independently of the real copy. Let
$d_{\rm rat}(b;t_I)$ be the density of the resulting true
integer polynomial output $b(t)\bmod N$, relative to uniform
measure on $(\Z/N\Z)^{m_{\rm out}}$. Every coefficient in
this congruence law is the actual integer coefficient of the
path. In particular, the small terms omitted from $G$ remain
in this law.

The comparison weight on the parent is
\begin{equation}
 K_*(u)=\kappa G(u/H,y_W,z_{\rm act}/K)
 \E_{t_I}\left[
   \left(\prod_{i\in I}K_i\right)
   1_{z_I(u)=z_I(t_I)}d_{\rm rat}(b(u);t_I)\right],
 \qquad
 \kappa=\prod_h
 \frac{\covol(\Lambda_{W_h})}
      {\prod_{i\ {\rm in\ layer}\ h}K_i}.
 \label{cube:forecast}
\end{equation}
Buffered chart cutoffs are included and equal one on the relevant
supports. All factors are nonnegative. The prefactor converts
Haar measure in the parent chart into real volume on the
continuous axes, normalized active grid volume, and exact
inactive counting measure. The exact equality on inactive
coordinates uses the same $t_I$ as the residue law.

The comparison required for each site-twist model atom $J$ is
\[
             \E_U J K_*\;\approx\;
                  \E_{t\ {\rm in\ the\ slice}}J(\psi(t)).
\]
We prove it in Proposition~\ref{cube:conditioned-transfer} by
testing one bounded modulus at a time. The large product $N$
defines a finite comparison law; it need not be an admissible
equidistribution modulus.

Two uniform bounds are needed before that comparison can be
used. The density $d_{\rm rat}$ must have a cap independent of
the number of primes, and $K_*$ must be approximable by
site-twist combinations with bounds independent of the eventual
prime cutoff. The real product-block estimate and its order-zero
integer version already control $G$ and the inactive law. The
remaining task is arithmetic: we must control the Fourier
coefficients of $d_{\rm rat}$ uniformly as more prime powers
are included.

\subsection{Modular rank and the bounds for the comparison weight}

For each degree tag, a primitive row of the leading homogeneous
polynomial can be differenced until its phase is linear. We
measure modular rank by the probability that this final linear
coefficient vanishes modulo a prime power. The next lemma turns
a small vanishing probability into Fourier decay. It allows
exceptional prime powers, provided that their product is bounded.
The following probabilistic argument establishes that product
bound for the coefficient law, with a bound chosen before the
prime cutoff. We then return to~\eqref{cube:forecast} to prove
its cap and site-twist approximation.

\subsubsection{A modular Fourier estimate}

Lemma~\ref{cube:modular-exceptions} and the coefficient-law
calculation following it will verify the hypotheses below for
the comparison weight just defined.

\begin{lemma}[Modular rank implies Fourier decay]\label{cube:modular-decay}
Fix $s\geq1$, $C>0$, integers $m_h\geq0$, and put $m=\sum_{h=1}^s m_h$.
For each parameter $z$, let
\[
 F_h(t;z)\in\mathbb Z[t_1,\ldots,t_v]^{m_h},\qquad \deg_t F_h\leq h.
\]
Suppose the homogeneous degree-$h$ part $H_h(t)$ is independent of $z$.
Let $N=\prod_{p\in\mathcal S}p^{A_p}$, where $\mathcal S$ is a finite set
of primes, and let $D\geq1$ be an integer.  Write $e_p=v_p(D)$.
Choose integers $0\leq b_p\leq A_p$, and put $D_0=\prod_{p\in\mathcal S}p^{b_p}$.
Suppose that for every $p\in\mathcal S$, every $b_p<a\leq A_p$, every
$h$ with $m_h>0$, and every row
$\lambda\in(\mathbb Z/p^a\mathbb Z)^{m_h}$ having a unit coordinate,
\begin{equation}\label{cube:modular-rank-hyp}
 \mathbb P_{u_1,\ldots,u_{h-1}}
 \left(\operatorname{Lin}_t
   \Delta_{u_1}\cdots\Delta_{u_{h-1}}(\lambda\cdot H_h)(t)=0
            \pmod {p^a}\right)\leq p^{-Ca}.
\end{equation}
Here the $u_i$ are independent and uniform in $(\mathbb Z/p^a\mathbb Z)^v$,
$\Delta_u P(t)=P(t+u)-P(t)$, and $\operatorname{Lin}_t$ means the vector
of coefficients of the linear terms.  For $h=1$ there are no differences.

Independently at each prime, sample $t\bmod p^{A_p}$ uniformly subject to
one prescribed coordinatewise congruence modulo $p^{\min(e_p,A_p)}$.
Combine these samples by the Chinese remainder theorem.  Let $\mu_z$ be the
law of $(F_h(t;z))_{h=1}^s$ on $(\mathbb Z/N\mathbb Z)^m$, and let $d_z$ be
its density relative to uniform measure.  Set
\[
 \gamma_s=\frac1{2^s s!},\qquad
 P=\frac{C\gamma_s}{2^s},\qquad
 K_s=\frac{4s}{\gamma_s},\qquad P_1=PK_s.
\]
For every character $\chi$ of exact order $q$,
\begin{equation}\label{cube:modular-fourier-decay}
 |\widehat d_z(\chi)|\leq (D_0D)^{P_1}q^{-P},
\end{equation}
uniformly in $z$ and the prescribed residues.
If $P>m+1$, deleting all Fourier terms of order greater than an integer
$T\geq1$ gives a uniform error at most
\begin{equation}\label{cube:modular-fourier-tail}
 \frac{(D_0D)^{P_1}}{P-m-1}\,T^{-(P-m-1)}.
\end{equation}
At a prime $p$ with $b_p=e_p=0$, the corresponding local density satisfies,
provided $P\geq m+2$,
\begin{equation}\label{cube:uncharged-prime-density}
 \|d_{z,p}-1\|_\infty\leq 2p^{-2}.
\end{equation}
\end{lemma}

\begin{proof}
We first work at one prime, suppressing its subscript.  A nontrivial
character has order $p^a$, with $1\leq a\leq A$.  Its restriction to the
$h$th group of coordinates has order $p^{a_h}$, where
$0\leq a_h\leq a$ and $\max_h a_h=a$.
An empty maximum of tag depths is interpreted as zero.
Choose initially an index $i$ with $a_i=a$.  Whenever a larger index $j>i$
has $a_j\geq a_i/(2i)$, replace $i$ by such a $j$.  This process terminates
at an index $h$ satisfying
\begin{equation}\label{cube:tag-gap}
 a_h\geq\gamma_s a,\qquad
 \max_{j>h}a_j<\frac{a_h}{2h}.
\end{equation}
Indeed the indices strictly increase, and the product of the possible
losses $2i$ is at most $2^{s-1}(s-1)!$.

Suppose first that $a>K_s(b+e)$.  Put
\[
 k=\max\{e,\max_{j>h}a_j\},\qquad B=a_h-hk.
\]
The definition of $K_s$ and \eqref{cube:tag-gap} imply
\[
 k<\frac{a_h}{2h},\qquad
 B>\frac{a_h}{2}\geq\frac{\gamma_s a}{2}>b.
\]
Partition the prescribed residue class further into classes modulo $p^k$.
In each class write $t=r+p^k x$, with $x$ uniform modulo $p^{a-k}$.
Every term from a tag $j>h$ is constant modulo its character denominator,
since integer polynomials respect congruences.

The remaining phase has degree at most $h$ in $x$.  Its homogeneous
part of degree $h$ is
\[
 \frac{\lambda\cdot H_h(x)}{p^B}\pmod1
\]
for a row $\lambda$ primitive modulo $p^B$.  This follows by expanding
$F_h(r+p^k x;z)$: its top part is $p^{hk}H_h(x)$, and the character on
tag $h$ has exact order $p^{a_h}$.  Terms of smaller degree may depend on
$r$ and $z$ and need no estimate.

Apply Cauchy--Schwarz $h-1$ times on the finite group
$(\mathbb Z/p^{a-k}\mathbb Z)^v$.  Averaging the resulting linear phase
in $x$ is zero unless its coefficient vector vanishes modulo $p^B$.
The difference variables project uniformly modulo $p^B$, because
$B\leq a-k$.  Thus \eqref{cube:modular-rank-hyp} gives, on every class,
\[
 |\mathbb E_x e(\text{phase}(x))|^{2^{h-1}}\leq p^{-CB},
 \qquad e(u)=\exp(2\pi i u).
\]
Consequently the original character average has modulus at most
\[
 p^{-CB/2^{h-1}}\leq p^{-Pa}.
\]
The same reasoning includes $h=1$, when the average of a linear
character is evaluated directly.

If instead $a\leq K_s(b+e)$, the trivial bound $1$ gives
\[
 |\widehat d_{z,p}(\chi)|
 \leq p^{P_1(b+e)}p^{-Pa}.
\]
This estimate therefore holds in both cases.  If $b=e=0$, only the first
case occurs, so there is no prefactor.  The local averages multiply by
the Chinese remainder theorem; their orders also multiply.  This proves
\eqref{cube:modular-fourier-decay}.

The number of characters of order exactly $q$ on
$(\mathbb Z/N\mathbb Z)^m$ is at most $q^m$: all such characters belong
to the subgroup killed by $q$, which has that size.  Hence the discarded
Fourier terms have total absolute mass at most
\[
 (D_0D)^{P_1}\sum_{q>T}q^{m-P}
 \leq \frac{(D_0D)^{P_1}}{P-m-1}\,T^{-(P-m-1)}.
\]
Finally, at an uncharged prime the nonconstant Fourier mass is at most
\[
 \sum_{a\geq1}p^{a(m-P)}
 =\frac{p^{m-P}}{1-p^{m-P}}\leq2p^{-2}.
\]
\end{proof}

\subsubsection{Producing the modular rank hypotheses}

We first give a probabilistic sufficient condition for
\eqref{cube:modular-rank-hyp}, then verify its coefficient-law
hypothesis using Corollary~\ref{sampling:finite-residues}.

\begin{lemma}[Independent coefficients and exceptional prime powers]
\label{cube:modular-exceptions}
Retain $s,m_h,m$ and $C$ from Lemma~\ref{cube:modular-decay}.  Put
\[
 c_s=2^{1-s},\qquad
 J\geq\left\lceil\frac{2(C+m+10)}{c_s}\right\rceil,
 \qquad Q_s=\max\{2,s^{4/c_s}\}.
\]
Consider random homogeneous vector polynomials $H_h$ modulo prime powers.
Assume that, for every component of $H_h$, there are $J$ squarefree
monomials of degree $h$, using pairwise disjoint sets of variables, whose
coefficients are independent uniform elements of the coefficient ring,
independently of all remaining coefficients.  Assume also that the
complete coefficient arrays at distinct primes are independent.

Call $p^a$ bad if \eqref{cube:modular-rank-hyp} fails for some tag and
primitive row, or if $p^a<Q_s$.  For any finite collection of tested
prime powers, including all lower depths, let $b_p$ be the largest bad
depth, with $b_p=0$ if none is bad.  Define the degree-dependent integer
\[
 B_s=\prod_p p^{\max(\{a\geq1:p^a<Q_s\}\cup\{0\})}.
\]
For every integer $R\geq2$,
\begin{equation}\label{cube:bad-prime-product}
 \mathbb P\left(\prod_p p^{b_p}>B_sR\right)
 \leq \frac{2}{9R^9}.
\end{equation}
This bound is independent of the number and depths of the tested primes.
\end{lemma}

\begin{proof}
We use a simple sublevel estimate.  If a multilinear polynomial of
degree at most $r$ has a unit coefficient on some squarefree monomial of
degree $r$, then, for uniform independent inputs modulo $p^a$,
\begin{equation}\label{cube:multilinear-sublevel}
 \mathbb P(p^a\mid f)\leq r p^{-a/2^{r-1}}\qquad(r\geq1).
\end{equation}
For $r=1$, condition on all but a variable having a unit coefficient.
For the induction, write $f=xg+h$ using a variable in the distinguished
monomial.  The polynomial $g$ has a unit coefficient on a monomial of
degree $r-1$.  Where $p^{\lceil a/2\rceil}\nmid g$, the conditional
probability of $p^a\mid xg+h$ is at most
$p^{\lceil a/2\rceil-1-a}\leq p^{-a/2}$.  On the complementary event
apply the induction hypothesis at depth $\lceil a/2\rceil$.  Adding the
two bounds proves \eqref{cube:multilinear-sublevel}.

Fix a primitive row on tag $h$ modulo $p^a$, and choose one of its unit
components.  Use the $J$ designated monomials of this component.  From
each monomial select one variable, and examine its linear coefficient
after $h-1$ differences.  The designated random coefficient from that
block occurs with multiplier equal to a permanent of an
$(h-1)\times(h-1)$ matrix of uniform difference variables.  No designated
coefficient from another block enters this linear coordinate, because
the variable sets of the designated monomials are disjoint.

Condition on the difference variables and all other polynomial
coefficients.  Each selected equation has conditional probability at
most
\[
 p^{-a}\gcd(p^a,\operatorname{perm}),
\]
and these bounds multiply.  The permanent has a unit coefficient on a
squarefree monomial of degree $h-1$.  Splitting according to whether its
valuation is at least $\lceil a/2\rceil$ and using
\eqref{cube:multilinear-sublevel} bounds the expected cost of one equation
by
\[
 h p^{-a/2^{h-1}}\leq s(p^a)^{-c_s}.
\]
For $h=1$ the multiplier is $1$ and the cost is exactly $p^{-a}$.
The shift variables used by the different permanents are disjoint, so
these expected bounds also multiply.  Thus, if $\rho$ denotes the
vanishing probability in \eqref{cube:modular-rank-hyp},
\[
 \mathbb E\rho\leq\bigl(s(p^a)^{-c_s}\bigr)^J.
\]
Markov's inequality and a union bound over at most $s(p^a)^m$ rows and
tags give
\[
 \mathbb P(p^a\text{ fails the rank test})
 \leq s^{J+1}(p^a)^{C+m-c_sJ}
 \leq(p^a)^{-10}\qquad(p^a\geq Q_s).
\]
The last inequality follows from the stated choices of $J,Q_s$.
For a product $d$ of powers of distinct primes, all at least $Q_s$,
independence across primes therefore bounds the probability that all its
chosen prime powers are bad by $d^{-10}$.

Put $W=\prod_{p:p^{b_p}\ge Q_s}p^{b_p}$, so that
$\prod_p p^{b_p}\le B_sW$.  This discards whole small factors;
it does not assert that dividing a bad prime power by a factor of
$B_s$ leaves a bad prime power.  If $W>R$, either a single factor
exceeds $R$, or a subcollection has product in $(R,R^2]$ (multiply them
one at a time until the product first exceeds $R$).  Union-bounding over
single-power witnesses and over these composite witnesses gives
\[
 2\sum_{d>R}d^{-10}\leq\frac{2}{9R^9}.
\]
The product of all largest bad powers is at most $B_s$ times the product
just considered, proving \eqref{cube:bad-prime-product}.
\end{proof}

\begin{remark}[Approximate coefficient laws]
Suppose every tested prime power is at most an integer $Q$, and the actual
joint coefficient law modulo each witness modulus is within $\eta$ in
total variation of a model satisfying Lemma~\ref{cube:modular-exceptions}.
It is enough to test single prime powers at most $Q$ and witness products
in $(R,R^2]$; no modulus exceeds $\max(Q,R^2)$.  The same union bound gives
\[
 \mathbb P\left(\prod_p p^{b_p}>B_sR\right)
 \leq\frac{2}{9R^9}+(Q+R^2)\eta.
\]
Consequently the bound on the exceptional-prime product can be chosen
before $Q$, provided the coefficient approximation accuracy can be
chosen after $Q$. For the sampling law,
Corollary~\ref{sampling:finite-residues} supplies this comparison
jointly with spatial residue conditions and the prescribed coefficient
tilt, as verified next.
\end{remark}


\paragraph{The probabilistic modular-rank requirement.}
Return to the integer output $b=(u,w,z_{\rm act})$ in the
comparison law, and let the slice stride be $D$.
The pure degree-$h$ part in the long variables is independent of
$t_I$ and of the translated origin of the slice.  It is tested
before dilation by $D$.

For a prime $p$, depth $a$, tag $h$, and primitive row on that tag
modulo $p^a$, apply $h-1$ independent complete differences to the
corresponding top homogeneous polynomial.  The desired good-depth
condition says that the probability its remaining linear
coefficient is zero modulo $p^a$ is at most $p^{-C_*a}$.
Write $m_{\rm out}=n+\sum_h(\dim W_h+\#\{\text{active axes in layer }h\})
\le n+\sum_h d_h$.  Choose the fixed-degree multiple $C_*$ so that
$C_*\gamma_s/2^s\ge m_{\rm out}+2$ uniformly over these choices.
For an integer cutoff $Q$, test precisely the depths $p^a\le Q$.
Declare a tested depth bad if this condition fails or $p^a<Q_s$,
and let $b_p$ be the largest such depth, or zero if none.
The required
probabilistic conclusion is
\begin{equation}
                  \prod_{p\le Q}p^{b_p}\le D_0
 \label{cube:bad-product}
\end{equation}
outside a prescribed small fraction of paths, with $D_0$ chosen
before $Q$.

In the ideal independent-residue model, each component has the $J$
disjoint monomials required by Lemma~\ref{cube:modular-exceptions}:
use dedicated blocks for active coordinates, free degree-$h$ coefficient
coordinates for $w$, and kernel linear coordinates for $u$.
Their selected coefficients are uniform and independent of all remaining
coefficients.  Choose $J$ by that Lemma with $C=C_*$, uniformly for
$m=m_{\rm out}\le n+\sum_h d_h$.  Thus the initial block count depends
only on the structural dimensions, before the active axes are classified.  Independence across primes then gives its
uniform exceptional-product tail.  Namely,
for a target exceptional fraction $\delta$, choose an integer
$R_{\rm exc}$ with $2/(9R_{\rm exc}^9)<\delta/2$ and put
$D_0=B_sR_{\rm exc}$.  The witnesses are single tested prime powers
and distinct-prime products in $(R_{\rm exc},R_{\rm exc}^2]$.
Their number is at most $Q+R_{\rm exc}^2$ and their moduli are at most
$M_{\max}=\max(Q,R_{\rm exc}^2)$.  No independence between different
depths at one prime is required.

Here is the quantitative passage to the actual tilted law.  For an
integer probability mass obtained by sampling a smooth normalized
interval density at scale $\ell$, with normalized bounded-variation
bound $B$, one has
\[
 \|\mathcal L(X\bmod M)-\operatorname{Unif}(\Z/M\Z)\|_{\rm TV}
       \le O(BM/\ell).
\]
Indeed each progression of spacing $M$ is a Riemann sum for
$1/M$ of the total mass.  Its error is controlled by the total
variation of the density; summing over the $M$ residues and
normalizing proves the bound.  For a product law the errors add
over the coordinates.

Corollary~\ref{sampling:finite-residues} compares the actual law, at
each witness modulus $M$, to the product coefficient-chart law jointly
with the unweighted spatial residue law.  On its Haar covers the
$w$ coefficients are exactly uniform deck coordinates, independent of
the projection coordinates and across coefficients.  The active
dedicated coefficients have interval lengths at least $c'A^h$.
On enormous axes all other pure-long degree-$h$ coefficient intervals
have lengths comparable to $\sigma K_i/L^h$; these tend to infinity
uniformly since $K_i>L^{C_s}$ and $C_s>h$.  Include these remaining
coefficients in the residue comparison as well.  Moderate
non-designated coefficients remain deterministic zeros.  Finally,
the spatial kernel coefficients have nondegenerate interval lengths
comparable, up to pre factors, to $H_i/L$, and every other random
spatial coefficient interval becomes broad at the later parent scales.
Thus the displayed residue estimate replaces all random entries of
the relevant homogeneous coefficient arrays by independent uniform
residues modulo $M$, keeping deterministic zeros as zeros.  CRT then
makes the complete arrays independent across the primes dividing $M$.

The number of scalar coefficients is a fixed polynomial of the
structural dimensions: the degree is fixed, the number of parameter
variables is structural, and there are only polynomially many
monomials.  Choose the total comparison error at each witness to be
at most $\eta=\delta/[2(Q+R_{\rm exc}^2)]$, dividing it between the
finite-residue corollary and the coordinatewise residue estimates.
Choose $A$ for the latter projection estimates after $Q$, and then
the rank and parent scales for the corollary and spatial estimates.
The logarithms of $M_{\max}$, the witness count, and $\eta^{-1}$
are permitted pre budgets.  The approximate-law remark now gives
an exceptional fraction below $\delta$.  In particular $D_0$ and all
structural block counts remain fixed before $Q$.

\subsubsection{Bounds for the comparison weight}

For paths satisfying~\eqref{cube:bad-product}, summing
Lemma~\ref{cube:modular-decay} with $P\ge m_{\rm out}+2$ gives a
pre cap $C_{\rm rat}$ for $d_{\rm rat}$ independent of the finite
prime collection.  Let $C_G$ bound $G$, and let $C_I$ be the cap
given by Lemma~\ref{cube:product-block} at cube order zero for the
joint inactive law relative to its grid volume.  Independence
across inactive axes gives that joint cap by multiplication; bounded
inactive axes are enumerated.  Positivity therefore gives
\[
 0\le K_*\le \kappa C_G C_{\rm rat}
       \left(\prod_{i\in I}K_i\right)\Pr(z_I(t_I)=z_I)
       \le\kappa C_G C_{\rm rat}C_I.
\]
Lemma~\ref{sampling:lattice-chart} bounds $\kappa$ independently of
subspace heights, so this is a pre cap.

Moreover $K_*$ is uniformly approximable, to any pre accuracy,
by a pre-cost sum of site twists.  Truncate the congruence Fourier
series to characters of individual order at most $T$, using its
uniform absolute tail.  For a character $\chi$ of order $q$, its
coefficient is
\[
 c_\chi(t_I)=\E_{t_\parallel}
                  \overline{\chi(b(t_\parallel,t_I))}.
\]
Since $b$ consists of integer polynomials, this depends only on
$t_I\bmod q$ and has modulus at most one.  Split into those at most
$q^{\dim t_I}$ coordinatewise residue classes, weighted by their
probabilities.  The order-zero inactive product-block estimate on
each class supplies the remaining smooth and residue site
characters.  If a conditioned axis is too short for that estimate,
its $K_i$ is pre-bounded in terms of $q$ and is enumerated.
Normalized projection coordinates are bounded linear functions of
the small lifts, while characters of $z,w$ reduce to characters of
$\beta$ modulo their individual moduli.  Thus multiplying by $G$
and $\chi(b)$ gives permitted site twists.  There are at most
$\sum_{q\le T}q^{m_{\rm out}}$ rational characters and the residue
counts have logarithms at most $(\dim t_I)\log T$.  Dividing the
desired error among these terms and using the product-block
coefficient bounds preserves pre cost.  No least common multiple
of all retained orders is required.

Enlarge the degree-zero seminorm by
$\sup_{K_*}|\E_U\phi K_*|$, allowing every finite prime collection
and every path satisfying its own bad-product condition.  This
definition does not assert that a random path satisfies infinitely
many tests.  The preceding uniform cap and approximation bounds
give detection for the enlarged seminorm by site twists alone:
if $|v|\le B$ and $|\E vK_*|\ge a$, approximate $K_*$ uniformly
to $a/(2B)$ by a twist combination of coefficient mass $M$; one
twist then correlates with $v$ by at least $a/(2M)$.
Lemma~\ref{cube:separation} still applies: the extra pointwise
envelope is just a pre-bounded constant.  All these model bounds can
therefore be fixed before choosing $Q$.

\subsection{Transferring the score and counting the new slots}


\begin{proof}[Proof of Proposition~\ref{cube:conditioned-transfer}]
We first fix uniform bounds for the approximation and its moduli.
We then compare the model terms with comparison weights and recover
a positive patch with the required number of slots.

We may assume $0<\epsilon\le1$.  First fix the exceptional-product
bound $D_0$ so that the modular failure probability will be less than
$\tau/4$, and use it in the enlarged comparison family.  Fix
$\eta\le\rho\tau/32$.  Lemma~\ref{cube:separation}, with the
uniform detection and comparison bounds above, supplies pre bounds
$N_0,M_0,b_1$ for the number of atoms, coefficient mass, and atom
complexity in an approximation at accuracy $\eta$.  These numerical
bounds are fixed before the ensemble scales or the actual atoms.

Choose a pre cutoff $Q$ containing every prime power that can divide
an atom modulus of complexity at most $b_1$.  Prescribe the
finite-residue errors and the model-comparison error
$\rho/[8(1+M_0)]$ using these uniform bounds.  Choose the later
parameters in the stated order; the simultaneous choices are detailed
in Section~\ref{cube:sampler-proof}.  The modular-rank estimate then
gives failure probability less than $\tau/4$ through this cutoff,
with the same $D_0$.

For this ensemble, apply the enlarged degree-zero approximation to
$h_0=(f-\lambda)T_*(u,\beta)$, extended by zero outside the buffered
charts.  Definition~\ref{counting:patch} gives $|h_0|\le1$.  We obtain
\[
 h_0=\sum_{i=1}^{N_0}c_iJ_i+r,\qquad
 \sum_i|c_i|\le M_0,
 \qquad X_0(r)\le\eta,
 \qquad \sup_{K_*}|\E_UrK_*|\le\eta,
\]
where $X_0$ is the degree-zero local-slice seminorm and the comparison
family includes every finite prime collection.  Markov's inequality gives
$\sup_B|\E_{t\in B}r(\psi(t))|\le\rho/8$ outside a set of paths
of probability at most $\tau/4$.  For a path satisfying
Equation~\eqref{cube:bad-product}, use the comparison weight for
\[
 N_Q=\prod_{p\le Q}p^{\lfloor\log Q/\log p\rfloor}.
\]
The comparisons below use only bounded individual covers, rather than
an equidistribution estimate modulo $N_Q$.

For each model twist one needs
\begin{equation}
 \E_U J_i K_*
       =\E_{t\ {\rm in\ the\ slice}}J_i(\psi(t))+o(1),
 \label{cube:model-comparison}
\end{equation}
with error at most $\rho/[8(1+M_0)]$.  Here are the individual
comparison steps.  Write an atom as
\[
 J_i(u)=F_i(u/H,y(u))\,R_i(u,\beta(u)\bmod q_i),
\]
where buffered cutoffs are included in $F_i$, and $F_i$ and $R_i$
have the permitted bounds.  A finite combination of such expressions
is handled by including its pre coefficient mass in the accuracy.
The integer maps between $\beta$ and $(z,w)$ convert the residue
factor to these coordinates without increasing its modulus.

Truncate the Fourier series of $d_{\rm rat}$ by individual order
at a pre threshold.  Lemma~\ref{cube:modular-decay} makes its
uniform absolute tail as small as needed; the cap for the averaged
inactive factor then bounds the resulting comparison error.  A
retained character of order $q$ and $J_i$ use only the common
modulus $\operatorname{lcm}(q,q_i)$.  For this term apply the
single-site Haar-cover equidistribution in
Corollary~\ref{sampling:normalization}, jointly with its spatial
residues and smooth spatial weight.  Its proof from
Lemma~\ref{sampling:analytic-removal} applies to these prescribed
Fourier approximations.  An exact inactive-grid condition is
extended in the buffered chart at width comparable to $1/K_j$;
because $K_j\le L^h$, its Fourier cost is an allowed late cost.
All other normalized chart maps have height-independent bounds.
Thus the error can be chosen uniformly in $t_I$ and in the fixed
path's coefficient ranges, and can be averaged over $t_I$.

For each layer the Haar sheet measure is Lebesgue measure in $y_W$
times counting measure in all $z_j$, divided by
$\operatorname{covol}(\Lambda_W)$.  Multiplying by
$\kappa\prod_{j\in I}K_j$ therefore leaves Lebesgue measure in
$y_W$, exact counting measure on $z_I$, and active grid volumes
$\prod_{j\notin I}K_j^{-1}$.  The spatial residues and the $w$
deck residues are independent uniform variables.  Sum an active
grid in each residue class required by this term and $J_i$.
Its normalized spacing is at most
$\operatorname{lcm}(q,q_i)/K_j$, and $K_j\ge(AL)^h$.
The pre Lipschitz bound for $G F_i$ makes these grid sums converge
uniformly to Lebesgue measure in $z_{\rm act}/K$, with independent
uniform residues.  This also verifies all normalization factors
in~\eqref{cube:forecast}.

Restore the full rational density using the same absolute tail.
The resulting integral samples real normalized $t_\parallel$ for
the principal smooth arguments, and independently samples its
allowed residue classes for the true integer polynomial output,
conditional on the same $t_I$.  Although that output density was
defined modulo $N_Q$, its integral against $R_i$ uses only its
marginal modulo $q_i$, since $q_i\mid N_Q$.  By the definition of
$d_{\rm rat}$ this marginal is exactly the true polynomial-output
law obtained from uniform allowed long-parameter residues modulo
$q_i$.

Finally partition each actual long side into those residue classes.
On each class $R_i$ is constant as a function of the parameter
residues; normalized Riemann summation for the smooth principal
arguments has error bounded by a pre constant times $q_i/L$.
The principal lift arguments differ from the true ones by
$O(\sigma)$ uniformly.  The constrained spatial law in Section~
\ref{sampling:constrained-section} gives the same bound for the total omitted
nonkernel spatial displacement, including the inactive columns.
Thus the actual and principal smooth test values differ by at most
a pre constant times $\sigma$.  All residue factors have used the
true integer path polynomials, so this comparison makes no claim
that small real errors preserve residues.  Choosing the truncation
accuracy, then $\sigma$, then $A$ and the equidistribution scales
gives the required error in~\eqref{cube:model-comparison}.

The productive paths have probability at least $\tau$, and the
two discarded sets each have probability at most $\tau/4$.
Choose a productive path outside both discarded sets and its
productive slice.
Equation~\eqref{cube:model-comparison} and the two residual bounds
give
\[
 \E_U(f-\lambda)T_*K_*
 \ge\rho-\rho/8-\rho/8-\eta\ge\rho/2=:S.
\]
Remove from $d_{\rm rat}$ the factors at primes that are neither
bad nor divisors of $D$ nor below a sufficiently large fixed cutoff.
The good-prime factors are pointwise $1+O(p^{-2})$, uniformly in
$t_I$.  Their product is within $1\pm\delta$ after the cutoff is
chosen in terms of $\epsilon$.  If $\widetilde K$ is the resulting
weight, then
\[
 (1-\delta)\widetilde K\le K_*\le(1+\delta)\widetilde K.
\]
Since $fT_*$ and $\lambda T_*$ are nonnegative,
\[
 \E(f-\lambda)T_*K_*
 \le(1+\delta)\E fT_*\widetilde K
       -(1-\delta)\lambda\E T_*\widetilde K.
\]
For $(1-\delta)/(1+\delta)\ge1-\epsilon$, it follows that
\[
 \E\bigl(f-(1-\epsilon)\lambda\bigr)T_*\widetilde K
 \ge S/(1+\delta).
\]
Thus the discount preserves an explicit score, without requiring
the deleted-prime error to be smaller than $S$.  If $P_\epsilon$
is the fixed small-prime cutoff, retained primes either have a bad
tested depth, divide $D$, or are at most $P_\epsilon$.  The product
of the largest bad powers is at most $D_0$, so there are at most
$\log D_0/\log2$ primes of the first kind.  Keeping their factors
to the largest tested depths uses one modulus $M_*$ with
\[
 \log M_*\le
 C(\log D_0+\log D+P_\epsilon)\log Q.
\]
This is a permitted polynomial logarithmic cost for fixed
$\epsilon$.

Partition the ambient box into allowed residue slices to fix
$u\bmod M_*$ and into fine normalized spatial boxes.  Freeze only
the smooth $u/H$ argument in $G$, leaving the polynomial arguments
of $T_*$ unchanged.  Its pre Lipschitz bound and the cap of the
remaining discrete factor allow a mesh for which the total score
error is at most $S/[2(1+\delta)]$.  Parent sides can be chosen
large enough that all nonempty partition pieces have the required
relative side bounds.  Averaging the normalized slice scores gives
one piece of score at least $S'=S/[2(1+\delta)]$.  On that piece
the remaining positive
weight is a smooth function of the small lifts, supplemented by
residue information about $\beta$ modulo $M_*$ and exact inactive
projection coordinates.

Finally put $r_h=\dim W_h$, choose an integral basis matrix $B_h$
of $W_h\cap\Z^{d_h}$ and a representative $c_h\in W_h$, and write
$C_h-c_h=B_hp_h(u)$.  The coordinates of $p_h$ have ordinary degree
at most $h$; their coefficients need not be bounded, by
Definition~\ref{counting:patch}.  This is the first point at which
orientation heights are charged.  Write $M=M_*$,
$b_h=\max(1,\|B_h\|_{\infty\to\infty})$, and
$q_h=\lceil16Mb_h\rceil$.  Use a smooth partition of the torus of
$p_h/M$ into $O(q_h^{r_h})$ cells, with local lifted supports of
radius $1/q_h$ and centers $a_h$.  Local lifts also handle cells
crossing a torus boundary.

On a cell write $p_h/M=k_h+r_h^{\rm loc}$.  The vector
$B_hMr_h^{\rm loc}-y_h=\beta_h-B_hMk_h$ is integral.
Across its contributing points its diameter in maximum norm is at
most $2Mb_h/q_h+1/2<1$.  It is therefore one fixed integral vector
$\beta_{h,0}$, giving
\begin{equation}
                     \beta_h=B_hM_*k_h+\beta_{h,0}
 \label{cube:slot-recovery}
\end{equation}
Introduce the $r_h$ new weight-$h$ residual coordinates
\[
 e_h=k_h-(p_h/M-a_h),\qquad
 y_h=B_hM(a_h-e_h)-\beta_{h,0}.
\]
Their supports lie in $[-1/4,1/4]^{r_h}$.  Since
$B_hMk_h\in W_h\cap M\Z^{d_h}$, all exact projection coordinates
and the residues of $\beta_h$ modulo $M$ are fixed in each cell.
The discrete factor of the recovered weight is consequently
constant there, and its remaining smooth factor is a function only
of the new residuals $e_h$.  Substitute
\eqref{cube:slot-recovery} into the original $d_*$-slot patch
formula, inserting the new determining slots before old slots of
equal weight.  This preserves every weighted degree and the
triangular order.  Multiply the old kernel by the cell cutoffs,
the positive recovered smooth weight, and the original buffered
chart cutoff evaluated at the recovered $y_h$.  Whenever this last
cutoff contributes, the displayed identity constructs a true small
integer lift; uniqueness identifies it with the original lift.
Thus the new formula has no false lifts.  Its kernel depends only
on the new and old residuals and has the required support.

Let $N_{\rm cell}$ be the number of product cells.  Then
\[
 \log N_{\rm cell}
 \le C\sum_hr_h(1+\log M+\log b_h).
\]
The normalized chart bounds and the affine formula for $y_h$ show
that the kernel's logarithmic Lipschitz bound is polynomial in
these same quantities and the previous budgets.  Its cap $C_*$ is
controlled as well.  Summing the nonnegative cell weights gives the
weight before this partition, so one cell piece has score at least
$S'/N_{\rm cell}$.  Divide its kernel by $\max(1,C_*)$.
The resulting positive patch has exactly
$d_*+\sum_h\dim W_h$ slots and
\[
 \log(\text{score}^{-1})
 \le\log(S'^{-1})+\log N_{\rm cell}+\log\max(1,C_*).
\]
This proves the claimed polynomial logarithmic score, complexity,
and slice bounds, with the permitted late height costs.
\end{proof}

\subsection{Completion of the constrained sampling theorem}
\label{cube:sampler-proof}

\begin{proof}[Proof of Proposition~\ref{sampling:constrained}]
We assemble the preceding estimates for one sampling law.  In particular,
the prime cutoffs used for transferring a fixed patch will be selected
from uniform model bounds before the law, rather than from a model
constructed for a different law.

\emph{Structural choices.}
Fix the degree bound $s$, the positive ambient dimension $n$, and the
dimensions $d_h$.  The cube orders required below are $0\le q\le s+1$;
order zero is used for comparison weights.  Choose the dedicated block
counts large enough for Lemma~\ref{cube:product-block} at every one of
these orders, including the fixed derivative orders used in
Proposition~\ref{cube:comparison}.  Choose the kernel and modular block
counts large enough for the real-minor and integer-cokernel estimates and
Lemma~\ref{cube:modular-exceptions}.  Each requirement is a fixed-degree
polynomial in $1+n+\sum_hd_h$, so their maximum is such a polynomial.
Choose the enormous-axis exponent larger than
\[
 \max_{1\le h\le s,\ 0\le q\le s+1}h\bigl(j(h,q)+1\bigr),
\]
and enlarge it by a fixed margin.  This dominates the power of $L$ in
the enormous-axis fiber error.  The resulting number $m$ of parameter
coordinates is now fixed.  None of these choices uses a later accuracy,
chart width, or testing budget.

\emph{Forecasting all pre bounds.}
Put the prescribed chart widths, slice costs, local test complexities,
input caps, and requested inverse accuracies into one pre budget.  The
constant-coefficient laws and all principal normalized coefficient widths
are then fixed as in Section~\ref{sampling:constrained-section}.  They have positive
pre widths; in particular they are not narrowed when $\sigma$ is chosen.
All inverse normalized domain volumes in the product-block estimates
are consequently pre-bounded, including on the allowed slices and their
bounded-modulus restrictions.

First forecast the detection constants and approximation bounds for
each $0\le r\le s$.  For transfer also fix a lower score $\rho>0$,
productive fraction $\tau>0$, and the prescribed fixed density discount.
Choose the exceptional-prime product bound $D_0$ so that the modular
exceptional fraction will be less than $\tau/4$.  Consider the family of
all comparison weights~\eqref{cube:forecast} with this bound $D_0$,
allowing every finite prime collection.  Its cap and its uniform
site-twist approximation bounds were proved above independently of
that collection and of the later scales.  Enlarge the degree-zero
local-slice seminorm by the supremum of its pairings with these weights.

The ordinary and enlarged seminorms satisfy the hypotheses of
Lemma~\ref{cube:separation} with pre data.  Indeed their local parts
are dominated by a pre multiple of the one-site density, and their
comparison parts have a pre-bounded pointwise envelope.  The one-site
cap threshold is fixed by the constant-coefficient cap and the full-scale
$x$ law.  Corollary~\ref{sampling:normalization} makes the mass above
that threshold as small as the separator requires by later scale choices.
Choose the separator's cap and detection threshold first, its required
tail error next, and then its coefficient-mass and number-of-terms bounds.
Thus for every bounded input the required approximation has at most
$N_0$ terms, coefficient mass at most $M_0$, and atom complexity at most
$b_1$, where
\[
 \log(2+N_0+M_0)+b_1\le\mathcal P(b_{\mathrm{pre}}).
\]
These are uniform numerical bounds, not a choice of the actual atoms.
The same bounds apply to the enlarged seminorm at accuracy
$\eta\le\rho\tau/32$.

Choose $Q$ from $b_1$ large enough to contain every prime power that
can divide any permitted atom modulus.  For example, an upper bound
for the modulus itself suffices.  Hence $\log Q$ has a pre bound.
Choose the modular comparison errors on the single-prime-power and
witness-product family after $Q$, as in the probabilistic modular-rank
argument.  That family has moduli at most
$\max(Q,R_{\mathrm{exc}}^2)$ and at most
$Q+R_{\mathrm{exc}}^2$ members; choosing the error per member at most
$\tau/[8(Q+R_{\mathrm{exc}}^2)]$ still costs only a pre logarithmic
budget.  Also prescribe the weight/model comparison error
$\rho/[8(1+M_0)]$, the Fourier truncation errors, and all forward,
noninjectivity, cap-tail, and cube errors.  Only a fixed number of
successive polynomial enlargements has been made.  Taking their maximum
allows all these requests to be met together.

\emph{One choice of later parameters.}
Choose $\sigma$ after the preceding pre requirements.  The joint
total-variation perturbation estimate in
Proposition~\ref{cube:comparison}, and the spatial bound
\eqref{sampling:nonkernel-support}, make the real replacement errors
smaller than those requirements.  Next choose $A$ large enough for all
product-block truncations, long-parameter residue Riemann sums,
enormous-axis lattice errors, and the projection-coefficient residue
comparisons through the already chosen cutoff $Q$.

Lemma~\ref{sampling:lattice-chart} and the gap selection in the
coefficient construction choose $L$ with
\[
 A\le L\le A^{C_s(1+\sum_hd_h)}
\]
for the actual spaces $W_h$.  All Fourier and equidistribution budgets
in Lemma~\ref{sampling:analytic-removal} and
Corollaries~\ref{sampling:normalization} and
\ref{sampling:finite-residues} have polynomial logarithms in the pre
budget, $\log\sigma^{-1}$, $\log L$, and the logarithmic accuracies,
individual moduli, and family counts just prescribed.  Replace
$\log L$ by $C_s(1+\sum_hd_h)\log A$ when choosing the rank threshold
$R$.  This makes that threshold valid throughout the allowed range,
without knowing which gap the spaces select.  Finally choose every
parent side sufficiently large for the resulting Fourier estimates,
integer spatial residue comparisons, fiber counting, boundary errors,
and noninjectivity estimate.  The narrow nonkernel intervals have widths
comparable to $\sigma H_i/(mT_j)$ and therefore also become sufficiently
long at this last stage.

For an explicit common majorant, choose $c_s\ge2$ so that $(2+x)^{c_s}$ dominates every
complete polynomial bound used at one stage below, including its
fixed multiplicative constant, and set
\begin{equation}\label{cube:budget-majorant}
 x_0=2+b_{\mathrm{pre}},\qquad
 x_{i+1}=(2+x_i)^{c_s}\quad(0\le i<14).
\end{equation}
The following table records the order of these bounds.  Dimensions,
structural coefficient-array counts, and complexities are bounded
directly. Fourier, model, and family counts, caps, and scales are
bounded through their logarithms; positive fractions, correlations,
and errors are measured by their inverse logarithms.
\begin{center}
\small
\begin{tabular}{@{}cp{0.85\linewidth}@{}}
\toprule
Bound & Quantities covered\\
\midrule
$x_1$ & Structural dimensions, coefficient counts, chart and slice data.\\
$x_2$ & $D_0$, comparison envelope, domination and one-site caps.\\
$x_3$ & Separator cap, detecting fraction, threshold, and slice mesh.\\
$x_4$ & Dual-Gowers output and the required averaged cube power.\\
$x_5$ & Cube-mixture and comparison-weight Fourier mass, count, and twist complexity.\\
$x_6$ & Inverse-theorem output complexity and detecting correlation.\\
$x_7$ & $N_0,M_0$, tail error, $Q$, and witness-family data.\\
$x_8$ & Model-comparison covers, truncations, and all final pre errors.\\
$x_9$ & $\sigma^{-1}$, using the minor, boundary, and perturbation bounds.\\
$x_{10}$ & $A$, using residue, Riemann-sum, and enormous-axis errors.\\
$x_{11}$ & $L$, uniformly over its allowed range.\\
$x_{12}$ & Late ambient Fourier budgets and individual covers.\\
$x_{13}$ & Rank threshold $R$.\\
$x_{14}$ & A sufficient lower bound for every parent side $H_i$.\\
\bottomrule
\end{tabular}
\end{center}
Here the product-axis Fourier calculation costs polynomially in the
number of axes: if $D$ axes each have coefficient mass at most $M$,
allocate error $\epsilon/[2D(M+1)^{D-1}]$ to each axis.  Its inverse
logarithm is
$\log(2D/\epsilon)+(D-1)\log(M+1)$, and logarithmic term counts
add.  The ambient Fourier approximation of an $a$-dimensional
Lipschitz function likewise costs $O(a\log(2K+1))$ for cutoff $K$;
it does not form a net of all Lipschitz functions.
The separator uses one detection call: in its notation, coefficient
mass $M$ is followed by at most
\[
 \left\lceil324C^3(1+C)M^2\eta^{-2}\right\rceil
\]
empirical terms.  Their number enters later budgets only through its
logarithm.  Each retained model comparison uses the product of two
individual moduli, never the product of all prime powers through $Q$.
Finally, the real perturbation uses fixed-size minors separately on
each output axis, so their sublevel exponents depend only on $s$.
These observations, together with the displayed estimates in the
preceding lemmas, justify every polynomial enlargement in the table.
There are fourteen enlargements, independently of all variable dimensions
and term counts.  Thus
\[
 \log\sigma^{-1},\ \log A,\ \log R,
 \ \text{and a sufficient }\log\min_iH_i
 \le x_{14}\le(2+b_{\mathrm{pre}})^{(2c_s)^{14}}.
\]
No bound involves a coefficient of the $C_h$ or a height of $W_h$.
The choices define the single law
\eqref{sampling:tilted-law} for each center tuple; its normalization is
$Z=1+o(1)>0$ uniformly in the centers by
Corollary~\ref{sampling:normalization}.

\emph{Property~\ref{sampling:property-geometry}.}
The spatial supports put every sampled box inside $U$.  On the support
of the coefficient density, the coefficientwise small lifts sum to a
polynomial staying in the site uniqueness chart.  The complementary
coefficients are ordinary integers.  Thus $y_h\circ\psi$ is fully slow
at the actual parameter scales and $\beta_h\circ\psi$ is an
integer-coefficient polynomial, both of degree at most $h$.
Corollary~\ref{sampling:normalization} supplies the prescribed
noninjectivity bound.  Its center-uniform normalization and constant-mode
cancellation also give the mixed-center one-site approximation uniformly
against unit-bounded inputs.  These assertions all concern the law just
chosen, with no conditioning on productive paths.

\emph{Property~\ref{sampling:property-detection}.}
Lemma~\ref{cube:pathwise-detection} applies to the law just chosen,
with all cube errors and mixture bounds prescribed above.  It gives
detection uniformly over inputs and path-dependent tests and slices,
including the degree-zero case.  Its proof uses a common meshed slice
and the nonnegative local cube powers, so the local niltests need not
be fixed across paths.

The one-site domination and tail error were already prescribed for
Lemma~\ref{cube:separation}.  Applying that lemma now gives the actual
finite approximations in
$X(\phi)=\E_\psi\sup_{B,F}|\E_{t\in B}\phi(\psi(t))F(t)|$,
with the forecast bounds $N_0,M_0,b_1$.  The atoms can depend on this
law and on $\phi$; their bounds do not.  No union bound over inputs or
over unrestricted polynomial-orbit coefficients is used.

\emph{Property~\ref{sampling:property-transfer}.}
For any fixed patch $T_*$ in the stated class and any $f,\lambda$ in
the stated ranges, the input $h_0=(f-\lambda)T_*$, extended by zero
outside the buffered charts, has cap one.  Apply the enlarged
degree-zero approximation to this input under the already fixed law.
Its actual atoms have modulus bounded by the previously fixed $b_1$,
so every prime power required by them is already covered by $Q$.
Their residual is small both on all but a $\tau/4$ fraction of local
paths and against every admissible comparison weight.  The modular
argument gives failure probability less than $\tau/4$ for this same
law, with the same $D_0$.  It is not a new sampler or a restriction
of the law defining the approximation seminorm.

The proof of Proposition~\ref{cube:conditioned-transfer} now applies
without another scale choice: its model comparison uses only the
individual covers for one retained rational Fourier character and
one model atom, whose bounds were forecast above.  In particular it
does not use equidistribution on the product of all prime powers
through $Q$. The proposition therefore gives the discounted positive
score with exactly $\sum_h\dim W_h$ new determining coordinates,
and with polynomial logarithmic complexity, inverse score, and
slice cost in the allowed augmented budget. This is precisely
Property~\ref{sampling:property-transfer}, including its
linear rank bound and its permitted late height dependence.

\emph{The unconstrained version.}
Set $\mathcal D=1$ and use the unconstrained spatial construction
in Section~\ref{sampling:constrained-section}.  For example, fix
$\ell_0(s)>(s+1)(s+3)$, enlarging it if needed by a degree-only
constant.  At order zero the root variable supplies the spatial fiber.
For $1\le q\le s+1$, a selected $q$-column
difference minor is a nonzero polynomial in a degree-bounded number
of normalized tuple variables.  The sublevel estimate controls its
small values.  The integer-character calculation for the independent
free columns gives the tail
$C(B'')^\ell\sum_{a>B}a^{q-\ell/q}$, which tends to zero with
$B$ because $\ell>q(q+2)$.  These bounds use no ambient spatial
dimension in the required number of columns.  On a good tuple the
root variable and these $q$ columns give the controlled spatial
pivot in each parent coordinate.  Taking the product of the
rowwise fiber bounds has only polynomial logarithmic cost in $n$.
The same finite-slice argument therefore proves cube comparison
and Property 2,
now without lift coordinates or a coefficient tilt.  They use only
the fixed comparability of the prescribed side lengths to $L$.
The nearly uniform $x$ law, small spatial displacements, and collision
bound give Property 1.  The smooth integer outer law of $(x,V)$ is
dominated by the required uniform product-box law by a pre-bounded
factor; the parameter average retains its prescribed uniform law.
Thus the same conclusions hold for distinct comparable prime sides
chosen in advance.  This proves every clause of the proposition.
\end{proof}
